\documentclass[final]{article}

\usepackage{aps360}
\usepackage[utf8]{inputenc}
\usepackage[T1]{fontenc}
\usepackage{hyperref}
\usepackage{url}
\usepackage{booktabs}
\usepackage{amsmath}
\usepackage{microtype}

\title{Sport-Informed Injury Risk Classification via Pose Estimation}
\author{%
  Harry Zhang\\
  1010988843, harryh.zhang@mail.utoronto.ca\\
  https://github.com/po8onthetrack/Deep-Learning-Project\\
}

\begin{document}

\maketitle
\vspace{-0.5in}

\section*{Introduction}
The project is motivated by the student's experience with badminton and bouldering injuries, including shoulder strain from poor smash mechanics and wrist sprain during dynamic bouldering moves. These examples show how many sport injuries are linked to measurable, preventable joint-angle deviations. The goal is to build a system that classifies body poses into injury-risk categories from a single image frame without motion-capture equipment. Deep learning is appropriate because fixed angle thresholds do not generalize across body sizes, camera viewpoints, or noisy keypoint estimates --- an MLP can learn a robust decision boundary from data where hard rules fail. Problem statement: given a 51-dimensional keypoint vector extracted from one image frame, classify the pose into one of four clinically grounded injury-risk classes.

\section*{Illustration}
\begin{figure}[!h]
\centering
\fbox{\begin{minipage}{0.93\linewidth}
\centering\small
\textbf{MPOSE2021 / LSP Extended images} $\rightarrow$
\textbf{Keypoint extraction (MoveNet)} $\rightarrow$
\textbf{51-dim vector: 17 joints $\times$ $(x,y,\mathrm{conf})$} $\rightarrow$
\textbf{REBA angle labelling} $\rightarrow$
\textbf{MLP(2 to 3layers, final layer outputs 4 categories} $\rightarrow$
\textbf{0 Neutral; 1 Trunk dev.\,$>30^\circ$; 2 Neck $>15^\circ$; 3 Overreach $>60^\circ$} $\rightarrow$
\textbf{Evaluate on held-out test set: Accuracy, Macro-F1, Per-class Precision/Recall}
\end{minipage}}
\caption{Training and evaluation pipeline. Keypoint vectors from MPOSE2021 (PoseNet, $K=17$) and LSP Extended (MoveNet) are auto-labelled via REBA-derived angle thresholds, then used to train and evaluate the MLP on held-out validation and test splits.}
\label{fig:pipeline}
\end{figure}

\section*{Background \& Related Work}
\textbf{REBA.} \citet{hignett2000reba} introduced a clinical postural scoring tool built from more than 600 postural examples. This project uses its trunk, neck, and upper-limb thresholds to generate biomechanically grounded labels instead of manual labels.

\textbf{BlazePose.} \citet{bazarevsky2020blazepose} showed that lightweight CNN pose estimation can run in real time on-device. This supports the choice of a frozen MoveNet-style backbone for fast keypoint extraction.

\textbf{HRNet.} \citet{sun2019hrnet} maintains high-resolution feature maps for accurate keypoint localization. My project trades some pose-estimation accuracy for speed by using a lighter frozen pose model.

\textbf{Deep REBA.} \citet{chatzis2022automatic} estimated REBA scores from RGB-derived 3D skeletons using a variational deep network. My project keeps the skeleton-plus-REBA idea but solves a simpler 2D, four-class problem.

\textbf{ConstructionPose3D.} \citet{fan2024cp3d} introduced a large construction-pose dataset for deep ergonomic assessment. My project extends pose-based risk classification to sports using consumer-camera data and no depth sensor. The gap is lightweight, sport-context injury-risk classification from monocular 2D pose with clinical labels.

\section*{Data Processing}
Two datasets are combined. MPOSE2021 \citep{mazzia2021action,mpose2021zenodo,pic4ser2021mpose} (PoseNet version, $K=17$, shape \texttt{(N,30,17,3)}) skews toward neutral everyday poses; temporal averaging collapses each 30-frame sequence to one 51-dim vector. Leeds Sports Pose Extended \citep{johnson2011lspet} adds 10,000 athlete images (parkour, gymnastics, athletics) rich in extreme joint configurations that map to Classes 1--3. MoveNet SinglePose Lightning extracts keypoints from LSP images in $(y,x,\mathrm{conf})$ format, swapped to $(x,y)$ to match MPOSE2021. Both sources share the 17-joint COCO ordering.

High-motion MPOSE2021 sequences are filtered by mean frame-to-frame displacement $s = \frac{1}{TJ}\sum_{t}\sum_{j}\|\mathbf{p}_{t+1,j}-\mathbf{p}_{t,j}\|_1$, retaining $\approx$6,000--8,000 samples. All original labels are discarded and replaced by REBA-derived angles $\theta = \arccos({\vec{BA}\cdot\vec{BC}}/{\|\vec{BA}\|\|\vec{BC}\|})$ for three triplets: neck (ear$\to$shoulder$\to$hip), trunk ($180^\circ$ minus shoulder$\to$hip$\to$knee), and overreach (hip$\to$shoulder$\to$elbow). The worst violation wins:

\begin{center}
\small
\begin{tabular}{lll}
\toprule
Class & Label & Condition \\
\midrule
0 & Neutral       & all joints within range \\
1 & Trunk flexion & trunk deviation $>30^\circ$ \\
2 & Neck strain   & neck angle $>15^\circ$ \\
3 & Overreach     & overreach angle $>60^\circ$ \\
\bottomrule
\end{tabular}
\end{center}

Thresholds are set at intermediate values relative to strict clinical REBA limits \citep{hignett2000reba} to ensure sufficient class representation. Low-confidence keypoints (conf\,$<0.2$) are excluded; each skeleton is centred at the hip and scaled by torso length. The combined dataset is split 70\,/\,15\,/\,15\% (train\,/\,validation\,/\,test) using stratified sampling; the test set is evaluated exactly once after all tuning is complete. Weighted cross-entropy with inverse class-frequency weights addresses class imbalance.

\section*{Architecture}
The MLP takes a 51-dimensional keypoint vector and outputs a four-class probability distribution: \texttt{Linear(51,128)} $\rightarrow$ ReLU $\rightarrow$ Dropout(0.3) $\rightarrow$ \texttt{Linear(128,64)} $\rightarrow$ ReLU $\rightarrow$ Dropout(0.3) $\rightarrow$ \texttt{Linear(64,4)} $\rightarrow$ Softmax. I will train with weighted cross-entropy loss, Adam (lr\,=\,0.001), 50 epochs, and early stopping with patience 10. An MLP fits this problem because the input is a structured numerical vector, not an image grid; a CNN would impose unnecessary spatial assumptions.

\section*{Baseline Model}
The baseline is a deterministic REBA-rule classifier with no learned parameters. It computes the same neck, trunk, upper-arm, and knee angles, applies the threshold table, and outputs the highest-risk class. This is a strong baseline because the training labels come from these rules. The MLP should only be considered useful if it improves accuracy and macro-F1 on the held-out test set by handling noisy keypoints, body-size variation, and camera-angle variation better than hard thresholds.

\section*{Ethical Considerations}
MPOSE2021 contains pose data derived from human videos, so I must confirm that its terms permit academic use. Although raw images are not used, the source datasets may not represent all body types, ages, or ethnicities, so errors may be uneven across users. False negatives could make an athlete overtrust the system; it should be framed as a training aid, not medical advice or a safety guarantee. Any sport or workplace deployment must obtain informed consent because pose tracking can become surveillance.

\newpage
\begin{thebibliography}{9}

\bibitem[Anthropic(n.d.a)]{anthropic2026claude}
Anthropic. (n.d.).
\newblock Claude.ai.
\newblock \url{https://claude.ai}. Accessed June 2, 2026.

\bibitem[Anthropic(n.d.b)]{anthropic2026claudecode}
Anthropic. (n.d.).
\newblock Claude Code.
\newblock \url{https://www.anthropic.com/claude-code}. Accessed June 2, 2026.

\bibitem[Bazarevsky et~al.(2020)]{bazarevsky2020blazepose}
Bazarevsky, V., Grishchenko, I., Raveendran, K., Zhu, T., Zhang, F., and Grundmann, M. (2020).
\newblock BlazePose: On-device real-time body pose tracking.
\newblock \textit{arXiv preprint arXiv:2006.10204}.

\bibitem[Fan et~al.(2024)]{fan2024cp3d}
Fan, C., Mei, Q., and Li, X. (2024).
\newblock 3D pose estimation dataset and deep learning-based ergonomic risk assessment in construction.
\newblock \textit{Automation in Construction}, 164, 105452.

\bibitem[Hignett and McAtamney(2000)]{hignett2000reba}
Hignett, S. and McAtamney, L. (2000).
\newblock Rapid entire body assessment (REBA).
\newblock \textit{Applied Ergonomics}, 31(2), 201--205.

\bibitem[Mazzia et~al.(2021)]{mazzia2021action}
Mazzia, V., Angarano, S., Salvetti, F., Angelini, F., and Chiaberge, M. (2021).
\newblock Action Transformer: A self-attention model for short-time pose-based human action recognition.
\newblock \textit{Pattern Recognition}, 124, 108487.

\bibitem[Mazzia et~al.(2021)]{mpose2021zenodo}
Mazzia, V., Angarano, S., Salvetti, F., Angelini, F., and Chiaberge, M. (2021).
\newblock MPOSE2021: A dataset for short-time pose-based human action recognition.
\newblock Zenodo. \url{https://zenodo.org/records/5803076}.

\bibitem[PIC4SeR(2021)]{pic4ser2021mpose}
PIC4SeR. (2021).
\newblock MPOSE2021\_Dataset.
\newblock \url{https://github.com/PIC4SeR/MPOSE2021_Dataset}.

\bibitem[Sun et~al.(2019)]{sun2019hrnet}
Sun, K., Xiao, B., Liu, D., and Wang, J. (2019).
\newblock Deep high-resolution representation learning for human pose estimation.
\newblock \textit{CVPR}, 5693--5703.

\bibitem[TensorFlow Hub(2024)]{tensorflowmovenet}
TensorFlow Hub. (2024).
\newblock MoveNet: Ultra fast and accurate pose detection model.
\newblock \url{https://www.tensorflow.org/hub/tutorials/movenet}.

\bibitem[Johnson and Everingham(2011)]{johnson2011lspet}
Johnson, S. and Everingham, M. (2011).
\newblock Learning effective human pose estimation from inaccurate annotation.
\newblock \textit{Proceedings of the IEEE Conference on Computer Vision and Pattern Recognition (CVPR)}, pp.\ 1465--1472.
\newblock \url{https://sam.johnson.io/research/lspet.html}.

\bibitem[Chatzis et~al.(2022)]{chatzis2022automatic}
Chatzis, T., Konstantinidis, D., and Dimitropoulos, K. (2022).
\newblock Automatic ergonomic risk assessment using a variational deep network architecture.
\newblock \textit{Sensors}, 22(16), 6051.

\end{thebibliography}

\end{document}
